\section{十五世纪的朝廷、法律与文化连结}

本节以朝廷程序、法律语言、家庭与性别、城市文化和海外联系为问题线索。正史、实录、律例与奏议通常更容易保存官员、男性精英和国家分类；女性、家庭成员、工匠、城市居民与海外中介则常仅在墓志、契约、地方志、碑刻、航海记录或案件中留下片段。材料本身的偏向，应当成为叙事的一部分。

\subsection{辅政、审案与制度程序}

\paragraph{1438（正统三年）十二月8日：陆川平绵之战：明讨伐勐茂司仁法，攻打邻近土司，但未能击败他} 此事把权力、法律或文化关系推进到可辨认的节点。账本的证据限度是编年候选的年序可由官修与后出材料交叉；具体月日、参与者、战果或数字必须回查所列材料，不能将单一叙事当作定论。可资核查的材料为《英宗实录》正统三年条；《明史·兵志》及相关列传；边镇、地方志或对方材料。它们能够说明制度如何表述自身、某种命令或交涉如何被记录，却未必能直接复原家庭内部关系、性别分工、城市日常或海外行动者的意图。事件发生的条件涉及前期边防、地方武装或跨区域资源调动的压力推动了该行动；其后留下的制度与社会后果为它改变了后续调兵、补给或地方秩序的条件；战果和伤亡须分开核对。译写候选以编年材料定位；实录记载反映朝廷选择，崇祯期须另以奏疏、地方及敌对政权材料交叉。

\paragraph{1438（正统三年）：本年朝廷围绕继承、官僚协调和政令传达形成连续记录，可用于核对宣德至正统的中枢运作。（材料核查重点：诏令或奏议的形成与发受链）} 此事把权力、法律或文化关系推进到可辨认的节点。账本的证据限度是来源导向的年度桥接条目：本年材料可定位相关政务线索；具体月日、发文机关、地域和执行结果仍须逐项回查，不能以制度文本或奏报替代实际效果。可资核查的材料为《英宗实录》正统三年条；《明史·本纪》及相关列传；诏令、奏疏或会典版本。它们能够说明制度如何表述自身、某种命令或交涉如何被记录，却未必能直接复原家庭内部关系、性别分工、城市日常或海外行动者的意图。事件发生的条件涉及继承、官僚协调或皇权运作的压力使问题进入朝廷程序；其后留下的制度与社会后果为它影响后续人事与政策议程；动机和派系标签不能由单一叙事裁定。桥接条目只标示可回查的年度问题与材料链；实录有中央选择性，崇祯期尤其需要奏疏、地方、朝鲜及满文材料交叉。；本条特别核对诏令或奏议的形成与发受链。

\paragraph{1438（正统三年）：本年朝廷围绕继承、官僚协调和政令传达形成连续记录，可用于核对宣德至正统的中枢运作。（材料核查重点：报称信息的来源、抄传与行政分类）} 此事把权力、法律或文化关系推进到可辨认的节点。账本的证据限度是来源导向的年度桥接条目：本年材料可定位相关政务线索；具体月日、发文机关、地域和执行结果仍须逐项回查，不能以制度文本或奏报替代实际效果。可资核查的材料为《英宗实录》正统三年条；《明史·本纪》及相关列传；诏令、奏疏或会典版本。它们能够说明制度如何表述自身、某种命令或交涉如何被记录，却未必能直接复原家庭内部关系、性别分工、城市日常或海外行动者的意图。事件发生的条件涉及继承、官僚协调或皇权运作的压力使问题进入朝廷程序；其后留下的制度与社会后果为它影响后续人事与政策议程；动机和派系标签不能由单一叙事裁定。桥接条目只标示可回查的年度问题与材料链；实录有中央选择性，崇祯期尤其需要奏疏、地方、朝鲜及满文材料交叉。；本条特别核对报称信息的来源、抄传与行政分类。

\paragraph{1439（正统四年）：华北平原及山东地区发生洪涝灾害} 此事把权力、法律或文化关系推进到可辨认的节点。账本的证据限度是编年候选的年序可由官修与后出材料交叉；具体月日、参与者、战果或数字必须回查所列材料，不能将单一叙事当作定论。可资核查的材料为《英宗实录》正统四年条；《明史·五行志》；受灾府县地方志与奏疏。它们能够说明制度如何表述自身、某种命令或交涉如何被记录，却未必能直接复原家庭内部关系、性别分工、城市日常或海外行动者的意图。事件发生的条件涉及自然风险与地方报告促成朝廷或地方的应对记录；其后留下的制度与社会后果为赈济、迁徙和损失的实际范围仍须以受灾地区材料分别核验。译写候选以编年材料定位；实录记载反映朝廷选择，崇祯期须另以奏疏、地方及敌对政权材料交叉。

\paragraph{1439（正统四年）：交趾、北边或辽东的军令与边报在本年构成军事决策线索，报称战果应与对方及地方材料比照。（材料核查重点：诏令或奏议的形成与发受链）} 此事把权力、法律或文化关系推进到可辨认的节点。账本的证据限度是来源导向的年度桥接条目：本年材料可定位相关政务线索；具体月日、发文机关、地域和执行结果仍须逐项回查，不能以制度文本或奏报替代实际效果。可资核查的材料为《英宗实录》正统四年条；《明史·兵志》及相关列传；边镇、地方志或对方材料。它们能够说明制度如何表述自身、某种命令或交涉如何被记录，却未必能直接复原家庭内部关系、性别分工、城市日常或海外行动者的意图。事件发生的条件涉及前期边防、地方武装或跨区域资源调动的压力推动了该行动；其后留下的制度与社会后果为它改变了后续调兵、补给或地方秩序的条件；战果和伤亡须分开核对。桥接条目只标示可回查的年度问题与材料链；实录有中央选择性，崇祯期尤其需要奏疏、地方、朝鲜及满文材料交叉。；本条特别核对诏令或奏议的形成与发受链。

\paragraph{1439（正统四年）：交趾、北边或辽东的军令与边报在本年构成军事决策线索，报称战果应与对方及地方材料比照。（材料核查重点：报称信息的来源、抄传与行政分类）} 此事把权力、法律或文化关系推进到可辨认的节点。账本的证据限度是来源导向的年度桥接条目：本年材料可定位相关政务线索；具体月日、发文机关、地域和执行结果仍须逐项回查，不能以制度文本或奏报替代实际效果。可资核查的材料为《英宗实录》正统四年条；《明史·兵志》及相关列传；边镇、地方志或对方材料。它们能够说明制度如何表述自身、某种命令或交涉如何被记录，却未必能直接复原家庭内部关系、性别分工、城市日常或海外行动者的意图。事件发生的条件涉及前期边防、地方武装或跨区域资源调动的压力推动了该行动；其后留下的制度与社会后果为它改变了后续调兵、补给或地方秩序的条件；战果和伤亡须分开核对。桥接条目只标示可回查的年度问题与材料链；实录有中央选择性，崇祯期尤其需要奏疏、地方、朝鲜及满文材料交叉。；本条特别核对报称信息的来源、抄传与行政分类。

\subsection{皇室家庭与继承礼制}

\paragraph{1439（正统四年）：交趾、北边或辽东的军令与边报在本年构成军事决策线索，报称战果应与对方及地方材料比照。（材料核查重点：同年地方材料所见的执行范围）} 此事把权力、法律或文化关系推进到可辨认的节点。账本的证据限度是来源导向的年度桥接条目：本年材料可定位相关政务线索；具体月日、发文机关、地域和执行结果仍须逐项回查，不能以制度文本或奏报替代实际效果。可资核查的材料为《英宗实录》正统四年条；《明史·兵志》及相关列传；边镇、地方志或对方材料。它们能够说明制度如何表述自身、某种命令或交涉如何被记录，却未必能直接复原家庭内部关系、性别分工、城市日常或海外行动者的意图。事件发生的条件涉及前期边防、地方武装或跨区域资源调动的压力推动了该行动；其后留下的制度与社会后果为它改变了后续调兵、补给或地方秩序的条件；战果和伤亡须分开核对。桥接条目只标示可回查的年度问题与材料链；实录有中央选择性，崇祯期尤其需要奏疏、地方、朝鲜及满文材料交叉。；本条特别核对同年地方材料所见的执行范围。

\paragraph{1440（正统五年）：苏州、江南、华北平原、山东等地遭受洪涝灾害} 此事把权力、法律或文化关系推进到可辨认的节点。账本的证据限度是编年候选的年序可由官修与后出材料交叉；具体月日、参与者、战果或数字必须回查所列材料，不能将单一叙事当作定论。可资核查的材料为《英宗实录》正统五年条；《明史·五行志》；受灾府县地方志与奏疏。它们能够说明制度如何表述自身、某种命令或交涉如何被记录，却未必能直接复原家庭内部关系、性别分工、城市日常或海外行动者的意图。事件发生的条件涉及自然风险与地方报告促成朝廷或地方的应对记录；其后留下的制度与社会后果为赈济、迁徙和损失的实际范围仍须以受灾地区材料分别核验。译写候选以编年材料定位；实录记载反映朝廷选择，崇祯期须另以奏疏、地方及敌对政权材料交叉。

\paragraph{1440（正统五年）：浙江饥荒袭来} 此事把权力、法律或文化关系推进到可辨认的节点。账本的证据限度是编年候选的年序可由官修与后出材料交叉；具体月日、参与者、战果或数字必须回查所列材料，不能将单一叙事当作定论。可资核查的材料为《英宗实录》正统五年条；《明史·五行志》；受灾府县地方志与奏疏。它们能够说明制度如何表述自身、某种命令或交涉如何被记录，却未必能直接复原家庭内部关系、性别分工、城市日常或海外行动者的意图。事件发生的条件涉及自然风险与地方报告促成朝廷或地方的应对记录；其后留下的制度与社会后果为赈济、迁徙和损失的实际范围仍须以受灾地区材料分别核验。译写候选以编年材料定位；实录记载反映朝廷选择，崇祯期须另以奏疏、地方及敌对政权材料交叉。

\paragraph{1440（正统五年）：漕运、粮饷与军户事务的年度文书显示财政—军事相互约束，定额不能替代实际征解} 此事把权力、法律或文化关系推进到可辨认的节点。账本的证据限度是来源导向的年度桥接条目：本年材料可定位相关政务线索；具体月日、发文机关、地域和执行结果仍须逐项回查，不能以制度文本或奏报替代实际效果。可资核查的材料为《英宗实录》正统五年条；《明会典》相关条目；地方志、黄册/契约/河工材料。它们能够说明制度如何表述自身、某种命令或交涉如何被记录，却未必能直接复原家庭内部关系、性别分工、城市日常或海外行动者的意图。事件发生的条件涉及财政征解、交通工程或货币支付的压力使相关措施进入政务；其后留下的制度与社会后果为其法定安排不等于地方执行，后续须以地域材料追踪负担与成效。桥接条目只标示可回查的年度问题与材料链；实录有中央选择性，崇祯期尤其需要奏疏、地方、朝鲜及满文材料交叉。

\paragraph{1441（正统六年）二月27日：陆川-平绵之战：明军进攻蒙茂} 此事把权力、法律或文化关系推进到可辨认的节点。账本的证据限度是编年候选的年序可由官修与后出材料交叉；具体月日、参与者、战果或数字必须回查所列材料，不能将单一叙事当作定论。可资核查的材料为《英宗实录》正统六年条；《明史·兵志》及相关列传；边镇、地方志或对方材料。它们能够说明制度如何表述自身、某种命令或交涉如何被记录，却未必能直接复原家庭内部关系、性别分工、城市日常或海外行动者的意图。事件发生的条件涉及前期边防、地方武装或跨区域资源调动的压力推动了该行动；其后留下的制度与社会后果为它改变了后续调兵、补给或地方秩序的条件；战果和伤亡须分开核对。译写候选以编年材料定位；实录记载反映朝廷选择，崇祯期须另以奏疏、地方及敌对政权材料交叉。

\paragraph{1441（正统六年）：华北平原及山东地区发生洪涝灾害} 此事把权力、法律或文化关系推进到可辨认的节点。账本的证据限度是编年候选的年序可由官修与后出材料交叉；具体月日、参与者、战果或数字必须回查所列材料，不能将单一叙事当作定论。可资核查的材料为《英宗实录》正统六年条；《明史·五行志》；受灾府县地方志与奏疏。它们能够说明制度如何表述自身、某种命令或交涉如何被记录，却未必能直接复原家庭内部关系、性别分工、城市日常或海外行动者的意图。事件发生的条件涉及自然风险与地方报告促成朝廷或地方的应对记录；其后留下的制度与社会后果为赈济、迁徙和损失的实际范围仍须以受灾地区材料分别核验。译写候选以编年材料定位；实录记载反映朝廷选择，崇祯期须另以奏疏、地方及敌对政权材料交叉。

\subsection{北京城市与文化空间}

\paragraph{1441（正统六年）：浙江饥荒袭来} 此事把权力、法律或文化关系推进到可辨认的节点。账本的证据限度是编年候选的年序可由官修与后出材料交叉；具体月日、参与者、战果或数字必须回查所列材料，不能将单一叙事当作定论。可资核查的材料为《英宗实录》正统六年条；《明史·五行志》；受灾府县地方志与奏疏。它们能够说明制度如何表述自身、某种命令或交涉如何被记录，却未必能直接复原家庭内部关系、性别分工、城市日常或海外行动者的意图。事件发生的条件涉及自然风险与地方报告促成朝廷或地方的应对记录；其后留下的制度与社会后果为赈济、迁徙和损失的实际范围仍须以受灾地区材料分别核验。译写候选以编年材料定位；实录记载反映朝廷选择，崇祯期须另以奏疏、地方及敌对政权材料交叉。

\paragraph{1442（正统七年）一月：陆川平绵之战：蒙茂战败，司仁法逃往阿瓦} 此事把权力、法律或文化关系推进到可辨认的节点。账本的证据限度是编年候选的年序可由官修与后出材料交叉；具体月日、参与者、战果或数字必须回查所列材料，不能将单一叙事当作定论。可资核查的材料为《英宗实录》正统七年条；《明史·兵志》及相关列传；边镇、地方志或对方材料。它们能够说明制度如何表述自身、某种命令或交涉如何被记录，却未必能直接复原家庭内部关系、性别分工、城市日常或海外行动者的意图。事件发生的条件涉及前期边防、地方武装或跨区域资源调动的压力推动了该行动；其后留下的制度与社会后果为它改变了后续调兵、补给或地方秩序的条件；战果和伤亡须分开核对。译写候选以编年材料定位；实录记载反映朝廷选择，崇祯期须另以奏疏、地方及敌对政权材料交叉。

\paragraph{1442（正统七年）十一月20日：张皇后（洪熙）去世} 此事把权力、法律或文化关系推进到可辨认的节点。账本的证据限度是编年候选的年序可由官修与后出材料交叉；具体月日、参与者、战果或数字必须回查所列材料，不能将单一叙事当作定论。可资核查的材料为《英宗实录》正统七年条；《明史·本纪》及相关列传；诏令、奏疏或会典版本。它们能够说明制度如何表述自身、某种命令或交涉如何被记录，却未必能直接复原家庭内部关系、性别分工、城市日常或海外行动者的意图。事件发生的条件涉及继承、官僚协调或皇权运作的压力使问题进入朝廷程序；其后留下的制度与社会后果为它影响后续人事与政策议程；动机和派系标签不能由单一叙事裁定。译写候选以编年材料定位；实录记载反映朝廷选择，崇祯期须另以奏疏、地方及敌对政权材料交叉。

\paragraph{1442（正统七年）：朝贡、册封和边地交涉的记录为本年多政权关系留档，礼仪语言与实际控制范围必须区分} 此事把权力、法律或文化关系推进到可辨认的节点。账本的证据限度是来源导向的年度桥接条目：本年材料可定位相关政务线索；具体月日、发文机关、地域和执行结果仍须逐项回查，不能以制度文本或奏报替代实际效果。可资核查的材料为《英宗实录》正统七年条；《明史·外国传》；《朝鲜王朝实录》、琉球《历代宝案》或海外档案。它们能够说明制度如何表述自身、某种命令或交涉如何被记录，却未必能直接复原家庭内部关系、性别分工、城市日常或海外行动者的意图。事件发生的条件涉及朝贡名义、边境安全与港口贸易的利益使交涉持续发生；其后留下的制度与社会后果为它为后续册封、互市或海防安排留下文书与跨政权比较线索。桥接条目只标示可回查的年度问题与材料链；实录有中央选择性，崇祯期尤其需要奏疏、地方、朝鲜及满文材料交叉。

\paragraph{1443（正统八年）三月：陆川平绵之战：明军击败司继发，但未能俘获他} 此事把权力、法律或文化关系推进到可辨认的节点。账本的证据限度是编年候选的年序可由官修与后出材料交叉；具体月日、参与者、战果或数字必须回查所列材料，不能将单一叙事当作定论。可资核查的材料为《英宗实录》正统八年条；《明史·兵志》及相关列传；边镇、地方志或对方材料。它们能够说明制度如何表述自身、某种命令或交涉如何被记录，却未必能直接复原家庭内部关系、性别分工、城市日常或海外行动者的意图。事件发生的条件涉及前期边防、地方武装或跨区域资源调动的压力推动了该行动；其后留下的制度与社会后果为它改变了后续调兵、补给或地方秩序的条件；战果和伤亡须分开核对。译写候选以编年材料定位；实录记载反映朝廷选择，崇祯期须另以奏疏、地方及敌对政权材料交叉。

\paragraph{1443（正统八年）：漕运、粮饷与军户事务的年度文书显示财政—军事相互约束，定额不能替代实际征解} 此事把权力、法律或文化关系推进到可辨认的节点。账本的证据限度是来源导向的年度桥接条目：本年材料可定位相关政务线索；具体月日、发文机关、地域和执行结果仍须逐项回查，不能以制度文本或奏报替代实际效果。可资核查的材料为《英宗实录》正统八年条；《明会典》相关条目；地方志、黄册/契约/河工材料。它们能够说明制度如何表述自身、某种命令或交涉如何被记录，却未必能直接复原家庭内部关系、性别分工、城市日常或海外行动者的意图。事件发生的条件涉及财政征解、交通工程或货币支付的压力使相关措施进入政务；其后留下的制度与社会后果为其法定安排不等于地方执行，后续须以地域材料追踪负担与成效。桥接条目只标示可回查的年度问题与材料链；实录有中央选择性，崇祯期尤其需要奏疏、地方、朝鲜及满文材料交叉。

\subsection{边疆交涉与外部观察}

\paragraph{1443（正统八年）：学校、科举和礼制的本年政策材料为社会流动和国家知识秩序提供可检索锚点} 此事把权力、法律或文化关系推进到可辨认的节点。账本的证据限度是来源导向的年度桥接条目：本年材料可定位相关政务线索；具体月日、发文机关、地域和执行结果仍须逐项回查，不能以制度文本或奏报替代实际效果。可资核查的材料为《英宗实录》正统八年条；《明史·选举志》或《艺文志》；文集、碑刻或书目材料。它们能够说明制度如何表述自身、某种命令或交涉如何被记录，却未必能直接复原家庭内部关系、性别分工、城市日常或海外行动者的意图。事件发生的条件涉及制度、教育或知识传播的需要推动了相关行动或文本形成；其后留下的制度与社会后果为它提供制度和传播史的锚点，不能据此推断社会整体接受。桥接条目只标示可回查的年度问题与材料链；实录有中央选择性，崇祯期尤其需要奏疏、地方、朝鲜及满文材料交叉。

\paragraph{1444（正统九年）：山西、陕西发生饥荒} 此事把权力、法律或文化关系推进到可辨认的节点。账本的证据限度是编年候选的年序可由官修与后出材料交叉；具体月日、参与者、战果或数字必须回查所列材料，不能将单一叙事当作定论。可资核查的材料为《英宗实录》正统九年条；《明史·五行志》；受灾府县地方志与奏疏。它们能够说明制度如何表述自身、某种命令或交涉如何被记录，却未必能直接复原家庭内部关系、性别分工、城市日常或海外行动者的意图。事件发生的条件涉及自然风险与地方报告促成朝廷或地方的应对记录；其后留下的制度与社会后果为赈济、迁徙和损失的实际范围仍须以受灾地区材料分别核验。译写候选以编年材料定位；实录记载反映朝廷选择，崇祯期须另以奏疏、地方及敌对政权材料交叉。

\paragraph{1444（正统九年）：苏北地区发生洪涝灾害} 此事把权力、法律或文化关系推进到可辨认的节点。账本的证据限度是编年候选的年序可由官修与后出材料交叉；具体月日、参与者、战果或数字必须回查所列材料，不能将单一叙事当作定论。可资核查的材料为《英宗实录》正统九年条；《明史·五行志》；受灾府县地方志与奏疏。它们能够说明制度如何表述自身、某种命令或交涉如何被记录，却未必能直接复原家庭内部关系、性别分工、城市日常或海外行动者的意图。事件发生的条件涉及自然风险与地方报告促成朝廷或地方的应对记录；其后留下的制度与社会后果为赈济、迁徙和损失的实际范围仍须以受灾地区材料分别核验。译写候选以编年材料定位；实录记载反映朝廷选择，崇祯期须另以奏疏、地方及敌对政权材料交叉。

\paragraph{1444（正统九年）：朝贡、册封和边地交涉的记录为本年多政权关系留档，礼仪语言与实际控制范围必须区分} 此事把权力、法律或文化关系推进到可辨认的节点。账本的证据限度是来源导向的年度桥接条目：本年材料可定位相关政务线索；具体月日、发文机关、地域和执行结果仍须逐项回查，不能以制度文本或奏报替代实际效果。可资核查的材料为《英宗实录》正统九年条；《明史·外国传》；《朝鲜王朝实录》、琉球《历代宝案》或海外档案。它们能够说明制度如何表述自身、某种命令或交涉如何被记录，却未必能直接复原家庭内部关系、性别分工、城市日常或海外行动者的意图。事件发生的条件涉及朝贡名义、边境安全与港口贸易的利益使交涉持续发生；其后留下的制度与社会后果为它为后续册封、互市或海防安排留下文书与跨政权比较线索。桥接条目只标示可回查的年度问题与材料链；实录有中央选择性，崇祯期尤其需要奏疏、地方、朝鲜及满文材料交叉。

\paragraph{1444（正统九年）：灾情上报、赈恤或河工材料标记本年地方风险，损失与救济效果需回到受灾地区核实} 此事把权力、法律或文化关系推进到可辨认的节点。账本的证据限度是来源导向的年度桥接条目：本年材料可定位相关政务线索；具体月日、发文机关、地域和执行结果仍须逐项回查，不能以制度文本或奏报替代实际效果。可资核查的材料为《英宗实录》正统九年条；《明史·五行志》；受灾府县地方志与奏疏。它们能够说明制度如何表述自身、某种命令或交涉如何被记录，却未必能直接复原家庭内部关系、性别分工、城市日常或海外行动者的意图。事件发生的条件涉及自然风险与地方报告促成朝廷或地方的应对记录；其后留下的制度与社会后果为赈济、迁徙和损失的实际范围仍须以受灾地区材料分别核验。桥接条目只标示可回查的年度问题与材料链；实录有中央选择性，崇祯期尤其需要奏疏、地方、朝鲜及满文材料交叉。

\paragraph{1445（正统十年）八月：陆川-平绵战役：阿瓦将斯仁法交给明朝，以换取明朝支持进攻鲜威} 此事把权力、法律或文化关系推进到可辨认的节点。账本的证据限度是编年候选的年序可由官修与后出材料交叉；具体月日、参与者、战果或数字必须回查所列材料，不能将单一叙事当作定论。可资核查的材料为《英宗实录》正统十年条；《明史·兵志》及相关列传；边镇、地方志或对方材料。它们能够说明制度如何表述自身、某种命令或交涉如何被记录，却未必能直接复原家庭内部关系、性别分工、城市日常或海外行动者的意图。事件发生的条件涉及前期边防、地方武装或跨区域资源调动的压力推动了该行动；其后留下的制度与社会后果为它改变了后续调兵、补给或地方秩序的条件；战果和伤亡须分开核对。译写候选以编年材料定位；实录记载反映朝廷选择，崇祯期须另以奏疏、地方及敌对政权材料交叉。
